\documentclass[12pt, % tamanho da fonte
	%openright,	 % capítulos começam em pág ímpar
	oneside,     %para impressão apenas em um lado (formato digital).	  
	%twoside,   %para impressão em frente e verso.  
	a4paper,     % tamanho do papel. 
	english,     % Idioma adicional para hifenização
	brazil       % Idioma principal 
	]{abntex2}
	
%------------------------------------------------------------
%------------    Estrutura do texto   -----------------------         

% Pacotes Básicos:
%\usepackage{lmodern}			    % Usa a fonte Latin Modern			

\usepackage[T1]{fontenc}		  % Selecao de codigos de fonte.
\usepackage[utf8]{inputenc}		% Codificacao do documento (conversão automática dos acentos)
\usepackage{mathptmx,helvet,courier}
%\usepackage{pslatex}                % Usa a fonte Times New Roman
\usepackage{lastpage}			    % Usado pela Ficha catalográfica
\usepackage[table]{xcolor}
\usepackage{color}				    % Controle das cores
\usepackage{graphicx}			    % Inclusão de gráficos
\usepackage{microtype} 		  	% para melhorias de justificação

\usepackage[normalem]{ulem}

\usepackage{quoting} % usar ambiente quoting (citações longas)

\usepackage{tablefootnote} % para colocar footnotes em tabelas e figuras
\usepackage{caption}
% Pacotes Extras:
\usepackage[final]{pdfpages}
\usepackage{amsmath,amsthm}   %Símbolos Matemáticos
%\usepackage{indentfirst} % Indenta primeiro parágrafo 
\usepackage[portuguese, ruled, linesnumbered,commentsnumbered, algo2e, vlined, lined, boxed, algochapter]{algorithm2e} % Algoritmos 
\usepackage{hyperref}
\usepackage[brazilian,hyperpageref]{backref}	 % Paginas com as citações na bibliograficas
% Configurações do pacote backref
% Usado sem a opção hyperpageref de backref
\renewcommand{\backrefpagesname}{Citado na(s) página(s):~}
% Texto padrão antes do número das páginas
\renewcommand{\backref}{}
% Define os textos da citação
\renewcommand*{\backrefalt}[4]{
\ifcase #1 %
Nenhuma citação no texto.%
\or
Citado na página #2.%
\else
Citado
\fi}%


%\usepackage[num]{abntex2cite}	% Citações padrão ABNT númerico
\usepackage[alf,abnt-emphasize=bf,abnt-full-initials=no,abnt-etal-list=5,abnt-etal-text=it,abnt-repeated-author-omit=yes]{abntex2cite}
%\usepackage[alf,abnt-emphasize=bf,abnt-etal-list=0,abnt-etal-text=emph]{abntex2cite}

% se desejar justificar as referências (pela abnt é alinhado a esquerda)
% \usepackage[alf,abnt-emphasize=bf,bibjustif]{abntex2cite}

\usepackage{setspace}

\usepackage{float} % para ajustar a posição das imagens

\usepackage{tikzsymbols} %para caracteres emoji
\usepackage{stackengine}
\usepackage{scalerel}


\newcommand\dangersign[1][2ex]{%
  \renewcommand\stacktype{L}%
  \scaleto{\stackon[1.3pt]{\color{red}$\triangle$}{\tiny\bfseries !}}{#1}%
}

\usepackage{etoolbox}
\usepackage{lipsum}
\usepackage{amssymb}
\usepackage{amsmath}

%para usar os algoritmos e matematica
\newcommand{\var}{\texttt}

% Defininfo Cores:
\definecolor{blue}{RGB}{25,25,112}
\definecolor{midgray}{gray}{.7}

\makeatletter % informações do PDF
\hypersetup{ % pagebackref=true,
	pdftitle={\@title}, 
	pdfauthor={\@author},
    pdfsubject={\imprimirpreambulo},
	pdfcreator={LaTeX with abnTeX2},
	pdfkeywords={abnt}{latex}{abntex}{abntex2}{trabalho acadêmico}, 
	colorlinks=true,     % false: boxed links; true: colored links
    linkcolor=black,          	% color of internal links
    citecolor=black,        		% color of links to bibliography
    filecolor=magenta,      	% color of file links
    urlcolor=black,
	bookmarksdepth=4 }
\makeatother
 

% -------------------------------------------- 
% Espaçamentos entre linhas e parágrafos 
\setlength{\parindent}{1.3cm} % O tamanho do parágrafo

% Controle do espaçamento entre um parágrafo e outro:
\setlength{\parskip}{0.2cm}  % tente também \onelineskip

% Definição de ambientes matemáticos em português 
\newtheorem{teorema}{Teorema}[chapter]
\newtheorem{axioma}{Axioma}[chapter]
\newtheorem{corolario}{Corolário}[chapter]
\newtheorem{lema}{Lema}[chapter]
\newtheorem{proposicao}{Proposição}[chapter]
\newtheorem{definicao}{Definição}[chapter]
\newtheorem{exemplo}{Exemplo}[chapter]
\newtheorem{observacao}{Observação}[chapter]

% Novos Comandos
\usepackage{tgtermes}

% Evita com que todas as seções fiquem com o estilo do título
\renewcommand{\tocheadstart}{\normalshape}

%%%%%%%%%%%%%%%%%%%%%%%%%%%%%%%%%%%%%%%%%%%%%%%%%%%%%%%%%%
% mostra no máximo subsubsection no sumário
\settocdepth{subsubsection} 

% diminuir espaço entre número e título do capítulo/seção/etc
\addtolength{\cftlastnumwidth}{-1.5em}
\cftsetindents{section}{0em}{\cftlastnumwidth}
\cftsetindents{subsection}{0em}{\cftlastnumwidth}
\cftsetindents{subsubsection}{0em}{\cftlastnumwidth}

% números de seções sem qualquer destaque
% seção primária (capítulo)
\let\oldthechapter\thechapter
\renewcommand{\thechapter}{\textnormal{\oldthechapter}}

% seção secundária (section)
\let\oldthesection\thesection
\renewcommand{\thesection}{\textnormal{\oldthesection}}

% seção terciária (subsection)
\let\oldthesubsection\thesubsection
\renewcommand{\thesubsection}{\textnormal{\oldthesubsection}}

% seção quaternária (subsubsection)
\let\oldthesubsubsection\thesubsubsection
\renewcommand{\thesubsubsection}{\textnormal{\oldthesubsubsection}}

% seção quinária (subsubsubsection)
\let\oldthesubsubsubsection\theparagraph
\renewcommand{\theparagraph}{\textnormal{\oldthesubsubsubsection}}


%%%%%%%%%%%%%%%%%%%%%%%%%%%%%%%%%%%%%%%%%%%%%%%%%%%%%%%%%%
% Estilos das seções

% seção primária (capítulo) em letra maiúscula, em negrito
\renewcommand{\ABNTEXchapterfont}{\rmfamily\bfseries}
\renewcommand{\ABNTEXchapterfontsize}{\normalsize}
\setboolean{ABNTEXupperchapter}{true}

% seção secundária em negrito 
\renewcommand{\ABNTEXsectionfont}{\rmfamily\bfseries}
\renewcommand{\ABNTEXsectionfontsize}{\normalsize}
\setboolean{ABNTEXuppersection}{false}

% seção terciária negrito e itálico
\renewcommand{\ABNTEXsubsectionfont}{\rmfamily\bfseries\itshape}
\renewcommand{\ABNTEXsubsectionfontsize}{\normalsize}
\setboolean{ABNTEXuppersubsection}{false}

% seção quaternária normal
\renewcommand{\ABNTEXsubsubsectionfont}{\rmfamily}
\renewcommand{\ABNTEXsubsubsectionfontsize}{\normalsize}
\setboolean{ABNTEXuppersubsubsection}{false}

% seção quinária em itálico
\renewcommand{\ABNTEXsubsubsubsectionfont}{\rmfamily\itshape}
\renewcommand{\ABNTEXsubsubsubsectionfontsize}{\normalsize}
\setboolean{ABNTEXuppersubsubsubsection}{false}

%%%%%%%%%%%%%%%%%%%%%%%%%%%%%%%%%%%%%%%%%%%%%%%%%%%%%%%
% Configuração das fontes no sumário

\renewcommand{\cftchapterfont}{\ABNTEXchapterfont}
\renewcommand{\cftchapterpagefont}{\normalsize}

\renewcommand{\cftsectionfont}{\ABNTEXsectionfont}
\renewcommand{\cftsectionpagefont}{\normalsize}

\renewcommand{\cftsubsectionfont}{\ABNTEXsubsectionfont}
\renewcommand{\cftsubsectionpagefont}{\normalsize}
  
\renewcommand{\cftsubsubsectionfont}{\ABNTEXsubsubsectionfont}
\renewcommand{\cftsubsubsectionpagefont}{\normalsize}

\renewcommand{\cftparagraphfont}{\ABNTEXsubsubsubsectionfont}
\renewcommand{\cftparagraphpagefont}{\normalsize}

% legenda de figuras com fonte reduzida
\let\oldlegend\legend
\renewcommand{\legend}[1]{\oldlegend{\ABNTEXfontereduzida{#1}}}

% Variáveis adicionais
\providecommand{\imprimirautorcite}{}
\newcommand{\autorcite}[1]{\renewcommand{\imprimirautorcite}{#1}} 
\providecommand{\imprimirsigla}{}
\newcommand{\sigla}[1]{\renewcommand{\imprimirsigla}{#1}}
\providecommand{\imprimiruf}{}
\newcommand{\uf}[1]{\renewcommand{\imprimiruf}{#1}}
\providecommand{\imprimircurso}{}
\newcommand{\curso}[1]{\renewcommand{\imprimircurso}{#1}}
\providecommand{\imprimirinstituto}{}
\newcommand{\instituto}[1]{\renewcommand{\imprimirinstituto}{#1}}
\providecommand{\imprimirdepartamento}{}
\newcommand{\departamento}[1]{\renewcommand{\imprimirdepartamento}{#1}}
\providecommand{\imprimirano}{}
\newcommand{\ano}[1]{\renewcommand{\imprimirano}{#1}}
\providecommand{\imprimirdia}{}
\newcommand{\dia}[1]{\renewcommand{\imprimirdia}{#1}}
\providecommand{\imprimirmes}{}
\newcommand{\mes}[1]{\renewcommand{\imprimirmes}{#1}}
\providecommand{\imprimirgrau}{}
\newcommand{\grau}[1]{\renewcommand{\imprimirgrau}{#1}}
\providecommand{\imprimirexaminadorum}{}
\newcommand{\examinadorum}[1]{
    \renewcommand{\imprimirexaminadorum}{#1}}
\providecommand{\imprimirexaminadordois}{}
\newcommand{\examinadordois}[1]{
    \renewcommand{\imprimirexaminadordois}{#1}}
\providecommand{\imprimirexaminadortres}{}
\newcommand{\examinadortres}[1]{
    \renewcommand{\imprimirexaminadortres}{#1}}
\providecommand{\imprimirexaminadorquatro}{}
\newcommand{\examinadorquatro}[1]{
    \renewcommand{\imprimirexaminadorquatro}{#1}}
\providecommand{\imprimirttorientador}{}
\newcommand{\ttorientador}[1]{
    \renewcommand{\imprimirttorientador}{#1}} 
\providecommand{\imprimirttcoorientador}{}
\newcommand{\ttcoorientador}[1]{
    \renewcommand{\imprimirttcoorientador}{#1}}
\providecommand{\imprimirttexaminadorum}{}
\newcommand{\ttexaminadorum}[1]{
    \renewcommand{\imprimirttexaminadorum}{#1}}
\providecommand{\imprimirttexaminadordois}{}
\newcommand{\ttexaminadordois}[1]{\renewcommand{
        \imprimirttexaminadordois}{#1}}
\providecommand{\imprimirttexaminadortres}{}
\newcommand{\ttexaminadortres}[1]{
    \renewcommand{\imprimirttexaminadortres}{#1}}
\providecommand{\imprimirttexaminadorquatro}{}
\newcommand{\ttexaminadorquatro}[1]{
    \renewcommand{\imprimirttexaminadorquatro}{#1}}

% Cria o comando \subtitulo. A norma define que o TITULO deve ser em caixa alta
% negrito, mas o subtitulo deve ser em caixa baixa.
\providecommand{\imprimirsubtitulo}{}
\newcommand{\subtitulo}[1]{\renewcommand{\imprimirsubtitulo}{#1}}

%----------------------------------------------------
\renewcommand{\imprimircapa}{  % Capa 
\begin{capa}
\center
\MakeTextUppercase{\imprimirinstituicao}\\ \MakeTextUppercase{\imprimirinstituto} \\
%\imprimirdepartamento \\
\MakeTextUppercase{\imprimircurso} \\              
\vspace{5cm}
\imprimirautor\\
\vspace{3.5cm}
\MakeTextUppercase{\textbf{\imprimirtitulo}}  \\
\MakeTextLowercase{\textbf{\imprimirsubtitulo}}
\vspace{2cm}
\vfill
\imprimirlocal~-~\imprimiruf \\ 
\imprimirano 
\end{capa}   } % Capa

%----------------------------------------------------
\renewcommand{\imprimirfolhaderosto}{% folha de rosto
    \begin{center}
    {{\MakeTextUppercase \imprimirautor}}  \\
		\vfill
		{\textbf{\MakeTextUppercase{\imprimirtitulo}}\\
        \MakeTextLowercase{\textbf{\imprimirsubtitulo}}}
    \end{center}
    \vfill 
    \begin{flushright} 
    \parbox{0.48\linewidth}{
	\imprimirtipotrabalho~ apresentado ao curso de Bacharelado em \imprimircurso~ da \imprimirinstituicao~ - \imprimirinstituto~ para a obtenção do título de \imprimirgrau. \\
	\vfill
	{\imprimirorientadorRotulo}~\imprimirorientador \\
	\vfill 
	{\imprimircoorientadorRotulo}~\imprimircoorientador}
   \end{flushright} 
	 \vfill
   \begin{center}
   {\large{\imprimirlocal~- \imprimiruf \\ \imprimirano}}
   \end{center} }  % folha de rosto

%----------------------------------------------------


\renewcommand{\listfigurename}{LISTA DE ILUSTRAÇÕES}
\pdfbookmark[0]{\listfigurename}{lof}

%================================================================================
% Pacotes de citações
%================================================================================

% ---
\makeatletter
\newcommand\thefontsize{Fonte tamanho: \f@size pt}
\makeatother



% para não deixar uma figura sozinha no centro de uma página,
% mas no topo da página.
\makeatletter
\setlength{\@fptop}{0pt}
\makeatother

\DeclareMathOperator*{\argmin}{arg\,min}

\let\theoldtextual\textual
\renewcommand{\textual}{
\theoldtextual
\pagenumbering{arabic} 
\setcounter{page}{\thesheetsequence-1}
}

\newcommand{\listafiguras}{
\listoffigures*   % Cria a Lista de Figuras
\cleardoublepage
}

% ---
\newcommand{\listatabelas}{
\pdfbookmark[0]{\listtablename}{lot}
\listoftables*  % Cria a lista de Tabelas
\cleardoublepage
}

\renewcommand{\listalgorithmcfname}{Lista de algoritmos}
\newcommand{\listaalgoritmos}{
\pdfbookmark[0]{\listalgorithmcfname}{lof}
\listofalgorithmes   % Cria a lista de Tabelas
\cleardoublepage
}

\newcommand{\sumario}{
\pdfbookmark[0]{\contentsname}{toc}
\tableofcontents*
\cleardoublepage
}

%--  % Estrutura do Texto e Pacotes Principais
\usepackage{listings} % listings e xcolor são responsáveis pela adição de segmentos de código destacados
\usepackage{xcolor}
\lstset{
  breaklines=true,            % Enable line breaking
 % breakatwhitespace=true,     % Break only at white space if possible
  postbreak=\mbox{\textcolor{red}{$\hookrightarrow$}\space}, % Symbol at line break
  frame=single,               % Draw a box around code
  basicstyle=\ttfamily\small, % Font style
  backgroundcolor=\color{white},
  columns=flexible            % Makes sure text fits nicely
  extendedchars=true,              % Suporte básico a caracteres acentuados
  literate={á}{{\'a}}1 {é}{{\'e}}1 {í}{{\'i}}1 {ó}{{\'o}}1 {ú}{{\'u}}1
  {Á}{{\'A}}1 {É}{{\'E}}1 {Í}{{\'I}}1 {Ó}{{\'O}}1 {Ú}{{\'U}}1
  {ã}{{\~a}}1 {ẽ}{{\~e}}1 {ĩ}{{\~i}}1 {õ}{{\~o}}1 {ũ}{{\~u}}1
  {Ã}{{\~A}}1 {Ẽ}{{\~E}}1 {Ĩ}{{\~I}}1 {Õ}{{\~O}}1 {Ũ}{{\~U}}1
  {â}{{\^a}}1 {ê}{{\^e}}1 {î}{{\^i}}1 {ô}{{\^o}}1 {û}{{\^u}}1
  {Â}{{\^A}}1 {Ê}{{\^E}}1 {Î}{{\^I}}1 {Ô}{{\^O}}1 {Û}{{\^U}}1
  {ç}{{\c{c}}}1 {Ç}{{\c{C}}}1
}

% -- Informações para Capa e Folha de Rosto a serem editadas
\titulo{Desenvolvimento de visualizações para o acompanhamento  de chamados de plataforma ITSM utilizando o Power BI}
\subtitulo{}
\autor{Ayrton Cavalieri de Almeida} 
\autorcite{SOBRENOME, Nome}
\local{Rio das Ostras} \uf{RJ}
\data{XX de XXXXX de 2026} \dia{XX} \mes{XX} \ano{2026} %deixar sem preencher antes do TCC
\orientador[Orientadora:]{Profa. Dra. Maise Dantas da Silva}
%\orientador[Orientadora:]{Profa. Dra. Nome completo}
\ttorientador{UFF} % Instituição do orientador
%\coorientador{Prof. Dr. Nome do Coorientador}
%\coorientador[Coorientadora]{Profa. Dra. Nome da Coorientadora}
%\ttcoorientador{UFF} % Instituição do Coorientador
\instituicao{Universidade Federal Fluminense} \sigla{UFF}
\instituto{Instituto de Ciência e Tecnologia}
\departamento{Departamento de Computação}
\curso{Ciência da Computação}	
\tipotrabalho{Trabalho de conclusão de curso}
\grau{Bacharel em Ciência da Computação}

%------Nomes dos examinadores.  
\examinadorum{Prof\textsuperscript{a}. Dr\textsuperscript{a}. Nome da professora} 
\ttexaminadorum{UFF}
\examinadordois{Prof. Dr. Nome do professor} 
\ttexaminadordois{UFF}
%\examinadortres{Prof. Dr. Membro da Banca  3}
%\ttexaminadortres{Universidade Federal de ... - UFXX}
%\examinadorquatro{Prof. Dr. Membro da Banca  4}
%\ttexaminadorquatro{Universidade Federal de ... - UFXX}

% ------------------------------------------------------
\begin{document} % Início do documento
\frenchspacing  % Retira espaço obsoleto entre as frases.
% ----------------------------------------------------------
% -- Elementos Pré-Textuais: -------------------------------
\imprimircapa  % Capa
\imprimirfolhaderosto 
% ---------------------------------------------------------------
% ----------------  Ficha Catalográfica  -------------------------
% ---------------------------------------------------------------
% Modelo de ficha catalográfica. Você deverá substituir esta
% folha na versão final da monografia por um pdf fornecido pela 
% biblioteca. Salve o modelo oficial como ficha_catalografica.pdf
% e use o comando abaixo para inseri-lo na versão final do texto.

%\begin{fichacatalografica}
%    \includepdf{ficha_catalografica.pdf}
%\end{fichacatalografica}



%% Modelo de Como fazer a Ficha Catalográfica:

\begin{fichacatalografica}
	\sffamily
	\vspace*{\fill}					% Posição vertical
	\begin{center}					% Minipage Centralizado
	\fbox{\begin{minipage}[c][8cm]{13.5cm}		% Largura
	
	Ficha catalográfica a ser elaborada pelo aluno.
% 	\small
% 	\imprimirautorcite.
% 	%Sobrenome, Nome do autor
	
% 	\hspace{0.5cm}  \\
% 	\imprimirtitulo  / \imprimirautor. --, \imprimirano-
	
% 	\hspace{0.5cm} \pageref{LastPage} p. 1 :il. (colors; grafs; tabs).\\
	
% 	\hspace{0.5cm} \imprimirorientadorRotulo~\imprimirorientador\\
	
% 	\hspace{0.5cm}
% 	\parbox[t]{\textwidth}{\imprimirtipotrabalho~--\\ \imprimirinstituicao. \\
% 	\imprimirinstituto. \\ ~\imprimirdepartamento.}\\
	
% 	\hspace{0.5cm}
% 		1. Palavra-chave1.
% 		2. Palavra-chave2.
% 		2. Palavra-chave3.
% 		I. \imprimirorientador.
% 		II. \imprimirinstituicao.
% 		III. Título 			
	\end{minipage}}
	\end{center}
\end{fichacatalografica}



% ---------------------------------------------------------------
% ----------------  Folha de aprovação  -------------------------
% ---------------------------------------------------------------
% Modelo de Folha de aprovação. Você deverá substituir esta
% folha na versão final da monografia por um pdf fornecido pelo  
% colegiado do seu curso. Salve o modelo oficial como 
% folhadeaprovacao_final.pdf e use o comando abaixo para
% inseri-lo na versão final do texto.

%\begin{fichacatalografica}
%    \includepdf{folhadeaprovacao_final.pdf}
%\end{fichacatalografica} Esta folha será 


\begin{folhadeaprovacao}

\begin{center}
    { \imprimirautor}  \\
       \vspace{2cm} 
       \textbf{\MakeTextUppercase{\imprimirtitulo}}\\
       \textbf{\MakeTextLowercase{\imprimirsubtitulo}}
    \end{center}
    \vspace{1.5cm}
    \begin{flushright} 
    \parbox{0.48\linewidth}{
		\imprimirtipotrabalho~ apresentado ao Curso de \imprimircurso~ da \imprimirinstituicao~ - \imprimirinstituto~ para a obtenção do título de \imprimirgrau. 
	}
   \end{flushright} 
  % \vspace{2cm}
   \vfill
 
   \begin{center}
   {
   Aprovado em:~\imprimirdia /\ \imprimirmes/\ \imprimirano.\\
   \vspace{1cm}
   BANCA EXAMINADORA
    \vspace{1cm}
    \vfill
          \rule{15cm}{.1pt} \\
      {\imprimirorientador}~-~{\imprimirttorientador}~(Orientadora) 
      \vfill
      \ifdefvoid{\imprimircoorientador}{}{
      \rule{15cm}{.1pt} \\
      \imprimircoorientador~-~\imprimirttcoorientador~(Coorientador) 
       \vfill
      }
      \rule{15cm}{.1pt} \\
      {\imprimirexaminadorum}~-~{\imprimirttexaminadorum} %\\ Examinador
        \vfill
        \ifdefvoid{\imprimirexaminadordois}{}{
        \rule{15cm}{.1pt} \\
        \imprimirexaminadordois~-~\imprimirttexaminadordois %\\ Examinador
	\vfill
        }
        \ifdefvoid{\imprimirexaminadortres}{}{
        \rule{15cm}{.1pt} \\
        \imprimirexaminadortres~-~\imprimirttexaminadortres %\\ Examinador
	\vfill
        }
        \ifdefvoid{\imprimirexaminadorquatro}{}{
        \rule{15cm}{.1pt} \\
        \imprimirexaminadorquatro~-~\imprimirttexaminadorquatro %\\ Examinador
        }
   }
   \vfill
   {\imprimirlocal~- \imprimiruf \\ \imprimirano}
   \end{center} 
\end{folhadeaprovacao}
% --- 
\begin{dedicatoria}
   \vspace*{\fill}
   \begin{flushright} 
        \parbox{0.48\linewidth}{
		Área da dedicatória. A dedicatória É opcional e deve ser digitada na parte inferior da folha e, de preferência, justificada em relação às margens laterais. Use esse recurso para dedicar o seu trabalho a todos que colaboraram com sua caminhada até aqui. Não utilize negrito, itálico ou outros modelos e tamanhos de fonte.
	}
   \end{flushright} 
   \vspace{2cm}
\end{dedicatoria}
\begin{agradecimentos}
Podemos já vislumbrar o modo pelo qual a constante divulgação das informações estende o alcance e a importância dos métodos utilizados na avaliação de resultados. 
Percebemos, cada vez mais, que o desenvolvimento contínuo de distintas formas de atuação garante a contribuição de um grupo importante na determinação das condições inegavelmente apropriadas.
O empenho em analisar o desafiador cenário globalizado cumpre um papel essencial na formulação do impacto na agilidade decisória.
O incentivo ao avanço tecnológico, assim como o desafiador cenário globalizado acarreta um processo de reformulação e modernização dos métodos utilizados na avaliação de resultados.

\end{agradecimentos}
\begin{epigrafe}
\vspace*{\fill}
\begin{flushright}
\parbox{0.7\linewidth}{
É claro que a determinação clara de objetivos maximiza as possibilidades por conta do investimento em reciclagem técnica.
		
\raggedleft\textit{Lero Lero.}
}
\end{flushright}
\vspace{2cm}
\end{epigrafe}
\chapter*{Declaração sobre o uso de Inteligência Artificial}
\addcontentsline{toc}{chapter}{Declaração sobre o uso de Inteligência Artificial}

Declaro que, durante a elaboração deste Trabalho de Conclusão de Curso, foi utilizada a ferramenta LanguageTool, disponibilizada pela Learneo, Inc., por meio de sua interface web, exclusivamente como recurso de apoio à revisão da linguagem, e o ChatGPT, disponibilizado pela Open AI para apoio à formatação técnica em \LaTeX.

As ferramentas foram empregadas para auxiliar na revisão da redação técnica, na adequação de trechos ao padrão acadêmico, na organização de estruturas em \LaTeX\ e no esclarecimento de dúvidas relacionadas à formatação do documento.

Não houve utilização da ferramenta para qualquer parte intelectual deste trabalho. Todo o conteúdo apresentado é de autoria do discente. %, sendo o uso da inteligência artificial restrito ao apoio técnico na escrita e na preparação do documento.
%--------------------------------------------------------------------------
%--------------------- Resumo em Português --------------------------------
%--------------------------------------------------------------------------

%\setlength{\absparsep}{18pt} % ajusta o espaçamento dos parágrafos do resumo
\begin{resumo}
Neste relatório técnico serão detalhadas as técnicas utilizadas para a obtenção de dados de plataforma ITSM, bem como as escolhas necessárias para a montagem do painel de visualização de dados. O trabalho foi realizado no âmbito do estágio de Análise de Sistemas na EDF Brasil e as atividades foram feitas em regime presencial. As responsabilidades durante o período de estágio concentraram-se em manutenção de sistemas utilizados pela empresa, fornecimento de dados através das soluções vigentes e montagem e análise de relatórios para as áreas clientes.
\textcolor{red}{(Ficou resumido demais. O ideal é o dobro desse tamanho. Mas é melhor deixarmos para pensar nele mais pro final da escrita.)}

 \vspace{\onelineskip}
 \noindent
 \textbf{Palavras-chave}: Análise de Dados; REST; \textit{Power BI}, Power Query, DAX.
\end{resumo}

%--------------------------------------------------------------------------
%--------------------- Resumo em Inglês --------------------------------
%--------------------------------------------------------------------------
\begin{resumo}[ABSTRACT]
\begin{otherlanguage*}{english}
The abstract text must go here and have the paragraph tab, on the first line,
single-spaced and must contain 150 to 500 words. It should not exceed one sheet.
Attention: it does not have paragraph indentation.In the current world, the execution of program points brings about a process of reformulation and modernization of the required financial and administrative conditions. We can already envision how the consensus on the need for qualification indirectly impacts the reassessment of the sectoral budget. In this way, the perception of difficulties is part of a process of managing the intended indexes.
The commitment to analyzing the constant dissemination of information challenges the ability to balance alternatives to orthodox solutions. Clearly, the commitment among teams has not yet convincingly demonstrated that it will participate in changing corporate paradigms. Nonetheless, the current structure of the organization prepares us to face atypical situations resulting from forms of action. Therefore, the challenging globalized scenario offers an interesting opportunity to verify development guidelines for the future. On the other hand, the new structural model advocated here presents trends towards approving the maintenance of normative conduct rules.

\vspace{\onelineskip}
\noindent 
\textbf{Keywords}: sleep apnea syndrome; placebo; polysomnography; \textit{Aedes aegypti}.
\end{otherlanguage*}
\end{resumo}
\listafiguras
\listatabelas
% ---------------------------------------------------
% ------ Lista de abreviaturas e siglas -------------
% ---------------------------------------------------
\begin{siglas}
\item[] 
  \item[API] \textit{Application Programming Interface}
  \item[BI] \textit{Business Intelligence}
  \item[DAX] \textit{Data Analysis Expressions}
  \item[DBA] \textit{Database Administrator}
  \item[ETL] \textit{Extract, Transform, Load}
  \item[HTTP] \textit{Hypertext Markup Language}
  \item[ITSM] \textit{Information Technology Service Management}
  \item[KPI] \textit{Key Performance Indicator}
  \item[MS] Microsoft
  \item[PBI] \textit{Power BI}
  \item[REST] \textit{Representational State Transfer}
  \item[SQL] \textit{Structured Query Language}
  \item[TI] Tecnologia da Informação
  \item[TO] Tecnologia de Operação
  \item[URL] \textit{Uniform Resource Locator}
  \item[UX] \textit{User Experience}
\end{siglas}

%% ---------------------------------------------------
% ----------- Lista de símbolos ---------------------
% ---------------------------------------------------

\begin{simbolos}
  \item[$ A \circ B $] Função composta de B em A
  \item[$ \underset{A_1, A_2 ... A_n}{argmin}$] Função que minimiza o resultado de outra função através da escolha dos parâmetros $A_1$,$A_2$... $A_n$
  \item[$ \left\|\; \right\|^2 $] Módulo de um vetor 
  \item[$ A\to B$] Transformação de A para B
  \item[$ \psi $] Letra grega minúscula psi
  \item[$ \phi$] Letra grega minúscula phi
  \item[$\sigma$] Letra grega minúscula sigma
  \item[$\Gamma$] Letra grega maiúscula gamma
  \item[$\gamma$] Letra grega minúscula gamma
  \item[$\mu$] Letra grega minúscula mu
  \item[$\nu$] Letra grega minúscula nu
  \item[$\mathbb {R}$] Conjunto dos Números Reais
  \item[$\int $] Integral
\end{simbolos}
\sumario
% ----------------------------------------------------------
% -- Capítulos do Trabalho: --------------------------------
\textual 
\chapter{Introdução}
\label{cap:Introdução}

A análise de dados permite o conhecimento, por exemplo, de tendências, gargalos e predição de problemas, entre outras informações úteis, auxiliando na tomada de decisões estratégicas em vários setores da economia.  O acesso a essas informações por parte de empresas e governos permite que tais decisões sejam melhor embasadas, reduzindo custos, melhorando a distribuição de mercadorias e serviços, e podendo, inclusive, até salvar vidas.

Atualmente, é comum que organizações possuam uma miríade de ativos de TI (Tecnologia da Informação) à disposição, como computadores, impressoras, equipamentos de rede e \textit{softwares}, por exemplo, utilizados pelos trabalhadores. Fazer a gestão de todos esses ativos é um problema que muitas empresas enfrentam ao longo de suas vidas. Entre outros pontos, é necessário saber quais são os usuários dos recursos, se determinado recurso continua ou não sendo útil, se seu prazo de vida útil já se esgotou, se algum equipamento está apresentando mais defeitos que a média, além de controlar possíveis alterações feitas pelo setor de TI no ambiente de usuário.  Para auxiliar nessa tarefa, utilizam-se ferramentas que permitem a análise e visualização de dados junto aos sistemas ITSM (\textit{Information Technology Service Management} - Serviço de Gerenciamento de Tecnologia da Informação), que permitem o controle de chamados e ativos pertencentes a uma organização. 

Com acesso total à fonte de dados dos sistemas ITSM, é possível alimentar ferramentas de análise, utilizando rotinas \textcolor{red}{(minha dúvida é em relação à concordância, porque, como está, é como se ETL fosse uma rotina específica, e não um padrão )} de tratamento de dados, como a ETL (\textit{Extract}, \textit{Load}, \textit{Transform} - Extrair, Transformar, Carregar), por exemplo, que permitem a manipulação e o tratamento do que está contido na base de dados. Após a obtenção e o tratamento, visualizações podem ser definidas para permitir que o gestor possa antecipar decisões e negociar com outros departamentos e organizações com os quais tenha contato.

Dessa forma, passa-se de um modelo de gestão reativo para uma abordagem proativa, otimizando o gerenciamento e reduzindo custos. Em vez de aguardar a falha estrutural de um equipamento para solicitar sua substituição ou manter licenças de \textit{software} ociosas, gerando despesas invisíveis, a gestão de TI passa a atuar fundamentada em evidências claras.

\textcolor{red}{Precisei reorganizar os 2 parágrafos seguintes porque as frases que apontam o que o trabalho apresenta estavam espalhadas. Dê uma olhada se ficou algo redundante.}

Nas seções seguintes, este relatório  detalha o trabalho realizado pelo autor durante seu estágio, com foco em Análise de Dados, na empresa Électricité de France (EDF), de novembro de 2022 a novembro de 2024. 
Nesse período, o autor atuou no setor de TI da Usina Termelétrica Norte Fluminense (UTENF), pertencente à EDF Brasil, onde teve contato com as áreas de Infraestrutura e Sistemas, tendo ficado um ano em cada uma. Na área de Infraestrutura, aprendeu sobre o funcionamento dos equipamentos de TI, bem como de TO (Tecnologia de Operação), participou da manutenção de sistemas críticos e realizou catalogação e manutenção de ativos da área. No período em Sistemas, aplicou os conhecimentos de Banco de Dados e de Estatística para desenvolver relatórios, fornecer acesso controlado às informações dentro do ambiente Microsoft e criar consultas, provendo \textit{views} específicas segundo a demanda das outras áreas.

É demonstrado neste trabalho como a modelagem dos indicadores-chave de desempenho (KPIs - \textit{Key Performance Indicators}) e a construção de painéis interativos transformam um vasto volume de dados brutos em inteligência acionável, garantindo a otimização contínua do ambiente tecnológico e o direcionamento assertivo dos investimentos da área. 
É apresentado também o processo de elaboração de um relatório para acompanhamento de chamados registrados no sistema ITSM da empresa e o processo de desenvolvimento utilizado pelo autor para definir os tratamentos de dados e a escolha de apresentação dos mesmos. São demonstradas as rotinas de ETL utilizando a ferramenta Power BI da Microsoft, bem como a confecção de relatórios para a obtenção de \textit{insights} e a previsão de tendências.

O trabalho está organizado da seguinte forma. 
Nesta seção são apresentados alguns detalhes sobre a empresa: sua atuação, histórico, missão e valores, além de informações sobre o estágio do autor. Na Seção 2 são introduzidos o ambiente de dados e as rotinas para tratamento e carregamento desses dados na plataforma Power BI. A Seção 3 detalha a confecção de um relatório para o departamento de TI, demonstrando a aplicação dos conceitos abordados, os elementos utilizados que são disponibilizados pela ferramenta Power BI, e os recursos de interação disponíveis. Por fim, na Seção 4, são apresentadas as conclusões.


\section{A EDF}


A Électricité de France, mais conhecida pela sigla EDF, é uma empresa estatal francesa do ramo de energia elétrica. Criada por decisão governamental em 8 de abril de 1946, a empresa  atualmente possui o monopólio de toda a grade de energia francesa sob seu ``guarda-chuva''. \cite{EDF_Historico}


A expansão da EDF se deu por investimentos do Plano Marshall e do governo francês durante o período conhecido como os ``Trinta Gloriosos'', o que a tornou global. Seu primeiro investimento em solo estrangeiro foi realizado em 1992. \cite{EDF_Historico}


Em 2004, o mercado de energia francês foi desregulamentado, fazendo com que parte da EDF fosse  disponibilizada na bolsa de valores durante o ano de 2005. Isso levou a empresa a adotar uma divisão societária de capital misto, mantendo o governo como acionista majoritário. \cite{EDF_Historico} 


Em 2022, a EDF foi reestatizada \cite{EDF_Reestatizacao}. Segundo o governo francês, ``a crise climática e a situação geopolítica requerem decisões firmes para garantir a independência da França e a soberania energética'' \cite{Sen_Frances}.


A Figura 1 apresenta a logo da EDF em 2025.
\begin{figure}[!ht]
	\centering
	\caption{Logo da EDF - 2025.}
	\includegraphics[width=0.25\linewidth]{figuras/edf_logo.jpg}
	\caption*{Fonte: \cite{Logo_EDF}.}
	\label{fig:placeholder Figura 1}
\end{figure}


\subsection{A EDF no Brasil}

A EDF atua no Brasil desde 2004 nos ramos de geração e de transmissão de energia.
Considerando a geração de energia, a empresa é dona de um vasto portfólio de ativos, possuindo 2,7GW de capacidade instalada, entre ativos solares, eólicos, hídricos e termelétricas. 
Esses ativos, que antes eram divididos entre as entidades EDF Brasil e EDF Renováveis, foram unificados em 2025, após um processo de reestruturação, formando a EDF Power Solutions \cite{EDF_PS_Unificacao,EDF_PS_Brasil}.
Já no setor de transmissão de energia, a EDF arrematou, por meio de sua então subsidiária UTE Norte Fluminense SA, 6km de linhas de transmissão, com capacidade de transmissão de 300MVA, servindo ao Complexo Porto do Açu e à Santo Amaro, ambos no estado do RJ \cite{Leilao_Transmissao}.


\subsection{Missão da EDF}
A missão da EDF é:

``Our raison d’être is to build a net zero energy future with electricity and innovative solutions and services, to help save the planet and drive wellbeing and economic development.''  \cite{EDF_PS_Missao}


Em tradução livre: 
Nossa razão de ser é construir um futuro com energia sem emissões de $\text{CO}_2$, utilizando eletricidade, soluções e serviços inovadores para ajudar a salvar o planeta e impulsionar o bem-estar e o desenvolvimento econômico.


Considerando sua missão, os valores da EDF se baseiam fortemente na promoção da transição energética, da inclusão, da diversidade social e de um ambiente de trabalho seguro e saudável.

\section{Sobre o estágio de TI}


Após a entrada na EDF, o primeiro dia do estagiário é dedicado ao \textit{onboarding} do time de Recursos Humanos (RH), com  reuniões introdutórias e vídeos instrucionais sobre as políticas da empresa. Caso o estagiário seja alocado em uma planta industrial, também é realizado um \textit{onboarding} de segurança junto ao time de QSMA (Qualidade, Meio Ambiente e Segurança). 


Em seguida, o estagiário é direcionado para um dos times de atuação dentro do departamento de TI, passando a integrar a equipe de Infraestrutura ou a de Sistemas.
O time de Infraestrutura tem como responsabilidades:  manter a rede da empresa segura e operacional, zelar pelos ativos de TI (sejam eles físicos ou licenças de \textit{softwares}), configurar e manter os \textit{data centers} e a infraestrutura de nuvem, disponibilizar Máquinas Virtuais com  serviços aos usuários, e prestar assistência ao time de Operações em assuntos que envolvam sistemas de TO. Já o time de Sistemas tem como funções: manter os \textit{softwares} corporativos operacionais, zelar pela integridade dos dados dos Bancos de Dados, realizar implementações e modificações dos processos dos departamentos mediante requisição, disponibilizar dados para criação de visualizações, executar modificações e manutenções nas visualizações dos departamentos, e controlar o acesso dos usuários às visualizações, garantindo que não haja acessos não autorizados.


Durante a atuação do autor, limitada à duração máxima do programa de 2 anos, o time de Sistemas era composto por um Especialista em Sistemas, dois Analistas terceirizados especializados no Sistema IBM MAXIMO, um DBA (\textit{Database Administrator}) terceirizado e um estagiário. Esse time respondia pelo gerenciamento dos sistemas responsáveis pelas operações da empresa e pelo controle de acesso aos dados fornecidos. 
Eram esperadas de seu estagiário: 
\begin{itemize}
	\item a criação  e a manutenção de relatórios para diversos departamentos;
	\item  a criação e a manutenção do \textit{gateway} de dados utilizado para acesso aos dados da empresa pelo serviço na nuvem do Power BI;
	\item a confecção de consultas SQL (com acessos controlados pelo Especialista)  para a disponibilização de fluxos de dados na ferramenta Power BI; e
	\item a criação de novos conectores no Power BI para permitir a obtenção de dados de fontes diversas.
\end{itemize} 
\chapter{Metodologias e Tecnologias Utilizadas}
\label{cap:relacionados}

Na área de Análise de Dados, conhecimentos de bancos de dados (como o SQL), estruturas de dados, lógica de programação, \textit{User Experience} (UX), requisições REST e estatística são aprofundados durante o estágio, sendo necessários para o exercício da profissão no futuro. Também serão utilizados conhecimentos de visualização de dados para melhor entregar os resultados esperados para a área cliente, junto à rotina ETL, usada para tratar os dados provenientes da fonte, que será explicada ao longo do capítulo.

Sobre as ferramentas empregadas, neste trabalho se utilizarão o TOPdesk, sistema de ITSM utilizado pela organização, de onde proverão os dados empregados e o PowerBI, uma aplicação da Microsoft especializada em tratamento, e visualização de dados. Detalhes sobre as ferramentas e onde são empregadas serão apresentados ao longo do capítulo.

\section{Tratamento de Dados}
O tratamento de dados é um passo importante para a disponibilização deles. É com ela que executa-se o saneamento dos dados entregues por uma fonte, corrigindo erros, executam-se cálculos prévios que poderão ser úteis a uma análise futura, filtram-se os dados que não interessam, dentre outras funções que podem ser utilizadas para entregar a informação requisitada. Para isso, utilizamos dois \textit{pipelines}, o ETL e o ELT, que serão detalhados a seguir. Ambas as rotinas permitem a consolidação de fontes de dados de diferentes origens.

\subsection{Pipeline ETL}
\label{subsec:ETL_Pipeline}
A rotina ETL realiza a \textit{Extração, Transformação e Carregamento} dos dados de uma base fornecida. Inicialmente, extraindo-os da origem, permitindo a disponibilidade para um motor de transformação, que aplicará operações, como filtragem de dados indesejados, criação de colunas calculadas que serão utilizadas a posteriori, correções e validações. Por fim, os dados são carregados na base de dados de destino, podendo ser concomitantemente de origem diversas origens.

Geralmente, as três fases podem ser executadas em paralelo, acelerando a disponibilização dos dados. Porém, este paralelismo deve ser executado observando-se a quantidade de recurso computacional disponível e a natureza do que está sendo carregado, a fim de não causar condições de corrida que podem sujar o resultado. \cite{Guia_de_ETL}

\subsection{Pipeline ELT}
Na rotina ELT, o carregamento e a transformação são feitas de forma invertida, com isso, rotinas de tratamento podem ser empregadas sobre o conjunto de dados já carregado no destino, removendo o motor de transformação do \textit{pipeline}, e entregando a responsabilidade das transformações para o destino. Neste caso, o destino precisa ter capacidade computacional necessária para realizar as transformações. \cite{Guia_de_ETL}

\subsection{ELT vs. ETL}
Os pipelines são empregados para cenários diferentes de processamento de dados, sendo assim, é necessário ter uma ponderação antes da implementação.

O ELT é recomendável para ambientes que possuem uma base de dados de destino com armazenamento e processamento escalonáveis, onde as rotinas podem se aproveitar dos recursos disponíveis e podendo se ajustar às necessidades dos usuários. Também recomenda-se ser utilizado caso queira preservar a base de dados original para realizar novas transformações no futuro, ou análise exploratória.

O ETL é recomendável quando a base de destino tem recursos limitados para a execução das transformações, podendo elas serem feitas com recurso computacional a parte. Também pode ser usado caso queira-se aplicar regras de negócio complexas, ou quando faz-se necessária a auditoria ou aplicação de regulamentações aos dados antes de serem carregados. \cite{Guia_de_ETL}

\section{Tecnologias empregadas}
Neste trabalho foram empregadas duas ferramentas. A primeira é o sistema ITSM utilizado pela companhia, o TOPdesk, que mantém os dados de chamados que serão utilizados. A segunda ferramenta é o PowerBI, plataforma da Microsoft que permite a aquisição de inteligência para a tomada de decisões pelos gestores das áreas. 

\subsection{TOPdesk}
O TOPdesk é uma plataforma que fornece serviços de apoio e gerência com foco nos profissionais de TI, RH (recursos humanos) e finanças, como gestão de ativos e ITSM, além de uma base de conhecimento que serve como uma enciclopédia, permitindo um acesso estruturado aos manuais, referências e procedimentos que foram submetidos. 


Essa plataforma fornece uma API REST\footnote{Uma API REST é ``uma Interface de Programação de Aplicativos (API) que segue os princípios de \textit{design} do estilo arquitetônico REST (sigla em inglês para transferência de estado representacional), utilizado para conectar sistemas hipermídia distribuídos'' \cite{REST_API_Definition}.}, que permite interações com os serviços prestados, como, por exemplo, consulta às bases de dados da plataforma para obtenção de ativos, chamados e operadores que, no caso deste trabalho, refere-se aos profissionais da TI e prestadores de serviços associados. Além disso, ela disponibiliza uma  documentação de apoio ao desenvolvedor para o desenvolvimento de aplicações que se conectem e interajam com o sistema. \cite{TOPdesk_API} 

\subsection{Power BI}
Segundo a Microsoft, 
\begin{quoting}[rightmargin=0cm,leftmargin=4cm]  \small
	O Power BI é a plataforma de análise de negócios da Microsoft que ajuda você a transformar dados em \textit{insights} acionáveis. Seja você um usuário de negócios, criador de relatórios ou desenvolvedor, o Power BI oferece ferramentas e serviços integrados para se conectar, visualizar e compartilhar dados em toda a sua organização. \cite{PBI_Sobre}
\end{quoting}


Portanto, o Power BI é uma ferramenta proprietária da Microsoft, voltada à análise de dados para negócios, que permite o tratamento e a visualização de dados de diversas fontes, como bancos de dados, arquivos JSON ou XML, \textit{dataframes}\footnote{Segundo a Databricks, um \textit{dataframe} é ``...uma estrutura de dados que organiza os dados em uma tabela bidimensional de linhas e colunas, como uma planilha'' \cite{Dataframe_Definition}.} processados por Python e R e dados hospedados na \textit{Web}, por exemplo. Entre as fontes de dados hospedados, destacam-se: o fluxo de dados hospedados no serviço de PBI na nuvem da Microsoft, os documentos hospedados no serviço de SharePoint (também da Microsoft) e os documentos disponibilizados por requisições REST  \cite{PBI_Fontes_Dados}.


A ferramenta disponibiliza duas linguagens para processamento de dados: a Power Query (ou Script M) \cite{PBI_M_Referencia} e a DAX \cite{PBI_DAX_Referencia}, que atuam em momentos diferentes na confecção dos relatórios. Além disso, ela também suporta a criação de visuais personalizados, utilizando JavaScript e R, e o processamento de dados, utilizando Python e R  \cite{PBI_Visuais_Customizados}.


Graças às ferramentas que o Power BI disponibiliza, é possível desenvolver rotinas ETL dentro da aplicação ou usar rotinas já escritas em Python e R. Porém, ao utilizar essa última, é necessária a configuração de um \textit{gateway} de dados\footnote{Um \textit{gateway} de dados é, segundo a Microsoft, ``...é um aplicativo cliente do Windows instalado localmente que atua como uma ponte entre suas fontes de dados locais e serviços na nuvem da Microsoft'' (tradução livre) \cite{Data_Gateway_Def}.}, pois a computação será realizada localmente nessa máquina \cite{PBI_Gateway_Python_R}.



\subsubsection{Script M}
Uma das linguagens disponibilizadas pela plataforma é o Script M. Seu paradigma é funcional, sensível à caixa e semelhante ao F\# \cite{PBI_M_Referencia}. Ela descreve as etapas de transformação e carregamento dos dados. Nesse papel, ela é responsável pela realização das transformações conforme explicadas na subseção \ref{subsec:ETL_Pipeline}. 


\subsubsection{DAX}
Outra linguagem disponibilizada é o DAX, sigla para \textit{Data Analysis Expressions} \cite{PBI_DAX_Referencia}. Ela é uma biblioteca de funções que permite a criação de rotinas de pós-processamento dos dados. Contém funções para a criação de medidas, colunas e tabelas calculadas e fórmulas para a função de garantia de segurança em nível de linha, limitando o acesso aos dados segundo o usuário \cite{PBI_DAX_Overview} \cite{PBI_RLS}.
\chapter{Implementação}
\label{cap:relacionados}

Neste capítulo, começaremos com as rotina de Extração, Transformação e Carregamento dos dados (ETL), iniciando pela extração ao abordar como fazer para realizarmos a coleta dessas informações do serviço externo. A seguir, abordaremos a transformação, onde fazemos o saneamento da base de dados baseado na amostra obtida na fase anterior, realizando o tratamento para que os dados estejam corretos, ou homogeneizados, o que é fundamental para o carregamento correto. No carregamento iremos efetivamente trazer os dados para a ferramenta e estabelecer as relações entre as tabelas de dados, utilizando as chaves que o desenvolvedor designou. Depois, veremos como fazer relatórios, utilizando as visualizações providas pela plataforma, que comunicam para o usuário informações que são relevantes e que auxiliam no processo de tomada de decisões estratégicas. Será abordado como ele influencia a tomada de decisões, permitindo que o gestor interaja com os mostradores, podendo obter detalhamentos e informações extras.

\section{ETL - \textit{Extract, Transform \& Load}}

Nesta sessão, discorreremos sobre a rotina ETL, sigla em inglês para Extração, Transformação e Carregamento. Ao finalizar essa rotina, os dados devem estar prontos para transformação, sendo possível montar relatórios e deles obter discernimento a cerca do que foi coletado da plataforma. Por fim, os dados são carregados na base de destino, podendo ser utilizado para a geração de relatórios.

\subsection{Extrair}

A Extração é a primeira etapa para obter os dados. Para isso, precisa-se de uma forma de conectarmos à base de dados e fazer o download dos registros necessários.

Como falado na introdução, o Power BI oferece muitos meios de conexão nativos, no caso do TOPdesk, temos uma situação onde um conector atende parcialmente a nossa necessidade, precisando de mais tratamentos para que a rotina seja executada. Para isso, utiliza-se a função nativa \texttt{Web.Contents}, que permite obter conteúdos de páginas web, configurar cabeçalhos, consultas, chaves de API, tempo de \textit{timeout}, e outros parâmetros retornando uma estrutura contendo a resposta do servidor, o conteúdo recebido e cabeçalhos de resposta \cite{WebContents}.

O TOPdesk retorna os dados da base em páginas com, no máximo, 10.000 linhas, não contemplando toda a base de dados sendo acessada, neste caso, o TOPdesk retornará em seu cabeçalho a resposta 206, e um endereço para o próximo segmento de dados no item \texttt{x-next-page-url}.\cite{TOPdesk_Incident_API} Porém, para reduzir o número de parâmetros, utiliza-se o número da página, passado na seção de consulta da URL, para navegar pelos dados. Esse conhecimento resulta na função \texttt{FetchData}, apresentada a seguir:

\begin{lstlisting}
FetchData = (pageNumber as number, pageSize as number, token as text, accumulated as list) =>
    let
        pagePosition = Number.ToText(pageSize * pageNumber),
	    header = [Authorization = Text.Combine({"Basic ", token})],
        source = Web.Contents("https://edf.topdesk.net/tas/api/incidents", [Query = [pageStart=pagePosition, pageSize=Number.ToText(pageSize)], Timeout=#duration(0, 0, 2, 0), Headers=header]),
        json = Json.Document(source),
        status = Value.Metadata(source)[Response.Status],
        result = if status = 206 then @FetchData(pageNumber + 1, pageSize, token, accumulated & {json}) else accumulated & {json}
    in
        result
\end{lstlisting}
\begin{center}
\text{Código da função FetchData}
\end{center}
Esta função recebe como parâmetros o número da página, a quantidade de linhas por página, a string com o token de autenticação e uma lista vazia para armazenar o conteúdo das respostas, em seguida executa os seguintes passos em ordem:
\begin{itemize}
    \item Calcula a posição da linha que inicia a página, para que a resposta desloque-se pelos dados e comece a entrega da linha requisitada, transformando o resultado numa string utilizando a função \texttt{Number.ToText}\cite{NumberToText};
    \item Monta o cabeçalho da requisição, gerando um \textit{registro} \cite{PBI_Tipos} que será passado para a função \texttt{Web.Contents} contendo o valor de "\textit{Authorization}", que indica para a página qual é o tipo de autenticação a ser utilizada e o token para autorizar. Neste caso, usa-se autenticação básica, passando um texto contendo a \textit{string} "\texttt{Basic }" concatenada com a \textit{string} de autenticação, que é formada pela concatenação do usuário e senha, codificados em base 64 \cite{TOPdesk_Login}. Em Script M, a concatenação de textos é feita utilizando a função nativa \texttt{Text.Combine}\cite{TextCombine};
    \item A função \texttt{Web.Contents} é chamada, passando como parâmetros o endereço do \textit{endpoint}\footnote{Segundo a Crowdstrike, um \textit{endpoint} é "...qualquer dispositivo físico que pode ser conectado a uma rede, incluindo computadores, laptops, celulares, tablets e servidores" \cite{Endpoint_Definition}.} da API do TOPdesk, um registro contendo a consulta (\textit{Query}) a ser concatenada na URL, sendo este parâmetro também um registro, contendo os valores \textit{pageStart} e \textit{pageSize} definidos anteriormente na função, o tempo máximo de espera (\textit{timeout}), sendo um parâmetro do tipo duração, e o cabeçalho, sendo este também um registro contendo o valor \textit{Authorization}. Recebendo como resposta um texto em formato JSON, o cabeçalho de resposta, com o código HTTP recebido, o conteúdo requisitado, e outros valores não utilizados neste contexto\cite{PBI_Tipos};
    \item A função \texttt{Json.Document} \cite{JsonDocument}, recebe a string em formato JSON obtida no passo anterior, retornando um registro com o conteúdo contido nela;
    \item Extrai os metadados associados à resposta de \texttt{Web.Contents}, utilizando a função \texttt{Value.Metadata} \cite{ValueMetadata}, obtendo o conteúdo indexado pela chave \texttt{Response.Status}.
\end{itemize}

Vê-se que a implementação é feita recursivamente, montando uma lista de listas, que na transformação serão concatenadas. Isso se deve ao fato de Script M ser uma linguagem funcional, não possuindo estrutura de loop como nas linguagens imperativas \cite{Functional_Loop}.

\subsection{Transformar}
A segunda etapa é a Transformação, nela realiza-se processamentos nos dados recebidos para que eles fiquem limpos e homogeneizados, permitindo o uso em construção de relatórios. As funções disponíveis no Script M permitem calcular colunas e tabelas, e corrigir erros encontrados no conjunto de dados obtido da plataforma ITSM.

\begin{lstlisting}
finalResult = List.Combine(FetchData(0, 10000, "chave_acesso", {})),
#"Converted to Table" = Table.FromList(finalResult, Splitter.SplitByNothing(), null, null, ExtraValues.Error),
#"Expanded Column1" = Table.ExpandRecordColumn(#"Converted to Table", "Column1", {"id", "status", "number", "request", "requests", "action", "attachments", "caller", "callerBranch", "callerLocation", "branchExtraFieldA", "branchExtraFieldB", "briefDescription", "externalNumber", "category", "subcategory", "callType", "entryType", "object", "asset", "branch", "location", "impact", "urgency", "priority", "duration", "actualDuration", "targetDate", "sla", "onHold", "onHoldDate", "onHoldDuration", "feedbackMessage", "feedbackRating", "operator", "operatorGroup", "supplier", "processingStatus", "responded", "responseDate", "completed", "completedDate", "closed", "closedDate", "closureCode", "timeSpent", "timeSpentFirstLine", "timeSpentSecondLine", "timeSpentPartial", "timeSpentLinkedPartials", "timeSpentSecondLineAndPartials", "itemCosts", "objectCosts", "costs", "escalationStatus", "escalationReason", "escalationOperator", "callDate", "creator", "creationDate", "modifier", "modificationDate", "majorCall", "majorCallObject", "publishToSsd", "monitored", "expectedTimeSpent", "archivingReason", "mainIncident", "optionalFields1", "optionalFields2", "externalLinks", "partialIncidents"}, {"id", "status", "number", "request", "requests", "action", "attachments", "caller", "callerBranch", "callerLocation", "branchExtraFieldA", "branchExtraFieldB", "briefDescription", "externalNumber", "category", "subcategory", "callType", "entryType", "object", "asset", "branch", "location", "impact", "urgency", "priority", "duration", "actualDuration", "targetDate", "sla", "onHold", "onHoldDate", "onHoldDuration", "feedbackMessage", "feedbackRating", "operator", "operatorGroup", "supplier", "processingStatus", "responded", "responseDate", "completed", "completedDate", "closed", "closedDate", "closureCode", "timeSpent", "timeSpentFirstLine", "timeSpentSecondLine", "timeSpentPartial", "timeSpentLinkedPartials", "timeSpentSecondLineAndPartials", "itemCosts", "objectCosts", "costs", "escalationStatus", "escalationReason", "escalationOperator", "callDate", "creator", "creationDate", "modifier", "modificationDate", "majorCall", "majorCallObject", "publishToSsd", "monitored", "expectedTimeSpent", "archivingReason", "mainIncident", "optionalFields1", "optionalFields2", "externalLinks", "partialIncidents"}),
#"Expanded caller" = Table.ExpandRecordColumn(#"Expanded Column1", "caller", {"id", "dynamicName", "email", "phoneNumber", "mobileNumber", "branch", "department", "budgetHolder"}, {"caller.id", "caller.dynamicName", "caller.email", "caller.phoneNumber", "caller.mobileNumber", "caller.branch", "caller.department", "caller.budgetHolder"}),
#"Expanded callerBranch" = Table.ExpandRecordColumn(#"Expanded caller", "callerBranch", {"name", "timeZone"}, {"callerBranch.name", "callerBranch.timeZone"}),
#"Expanded caller.department" = Table.ExpandRecordColumn(#"Expanded callerBranch", "caller.department", {"id", "name"}, {"caller.department.id", "caller.department.name"}),
#"Expanded category" = Table.ExpandRecordColumn(#"Expanded caller.department", "category", {"id", "name"}, {"category.id", "category.name"}),
#"Expanded subcategory" = Table.ExpandRecordColumn(#"Expanded category", "subcategory", {"id", "name"}, {"subcategory.id", "subcategory.name"}),
#"Expanded callType" = Table.ExpandRecordColumn(#"Expanded subcategory", "callType", {"name"}, {"callType.name"}),
#"Expanded entryType" = Table.ExpandRecordColumn(#"Expanded callType", "entryType", {"name"}, {"entryType.name"}),
#"Expanded asset" = Table.ExpandRecordColumn(#"Expanded entryType", "asset", {"id"}, {"asset.id"}),
#"Expanded priority" = Table.ExpandRecordColumn(#"Expanded asset", "priority", {"name"}, {"priority.name"}),
#"Expanded duration" = Table.ExpandRecordColumn(#"Expanded priority", "duration", {"name"}, {"duration.name"}),
#"Expanded sla" = Table.ExpandRecordColumn(#"Expanded duration", "sla", {"responseTargetDate", "targetDate"}, {"sla.responseTargetDate", "sla.targetDate"}),
#"Expanded operator" = Table.ExpandRecordColumn(#"Expanded sla", "operator", {"status", "name"}, {"operator.status", "operator.name"}),
#"Expanded operatorGroup" = Table.ExpandRecordColumn(#"Expanded operator", "operatorGroup", {"name"}, {"operatorGroup.name"}),
#"Duplicated Column" = Table.DuplicateColumn(#"Expanded operatorGroup", "operatorGroup.name", "operatorGroup.name - Original"),
#"Expanded processingStatus" = Table.ExpandRecordColumn(#"Duplicated Column", "processingStatus", {"name"}, {"processingStatus.name"}),
#"Expanded modifier" = Table.ExpandRecordColumn(#"Expanded processingStatus", "modifier", {"name"}, {"modifier.name"}),
#"Corrigindo RH" = Table.ReplaceValue(#"Expanded modifier","RH","Recursos Humanos",Replacer.ReplaceText,{"caller.department.name"}),
#"Corrigindo Elétrica" = Table.ReplaceValue(#"Corrigindo RH","Elétrica","Elétrico",Replacer.ReplaceText,{"caller.department.name"}),
#"Corrigindo C&L" = Table.ReplaceValue(#"Corrigindo Elétrica","Compras e Logistica","Compras e Logística",Replacer.ReplaceText,{"caller.department.name"}),
#"Corrigindo Mecânica" = Table.ReplaceValue(#"Corrigindo C&L","Mecânico","Mecânica",Replacer.ReplaceText,{"caller.department.name"}),
#"Corrigindo Mecânica1" = Table.ReplaceValue(#"Corrigindo Mecânica","Mecânida","Mecânica",Replacer.ReplaceText,{"caller.department.name"}),
#"Corrigindo Química" = Table.ReplaceValue(#"Corrigindo Mecânica1","Quimica","Química",Replacer.ReplaceText,{"caller.department.name"}),
#"Corrigindo TI" = Table.ReplaceValue(#"Corrigindo Química","Tecnologia da Informação","TI",Replacer.ReplaceText,{"caller.department.name"}),
#"Corrigindo TI1" = Table.ReplaceValue(#"Corrigindo TI","Tecnologia de Informação","TI",Replacer.ReplaceText,{"caller.department.name"}),
#"Corrigindo P&D" = Table.ReplaceValue(#"Corrigindo TI1","P&D","Pesquisa e Desenvolvimento",Replacer.ReplaceText,{"caller.department.name"}),
#"Corrigindo Marlim Azul" = Table.ReplaceValue(#"Corrigindo P&D","Marlim Azu","Marlim Azul",Replacer.ReplaceText,{"callerBranch.name"}),
#"Corrigindo Marlim Azul1" = Table.ReplaceValue(#"Corrigindo Marlim Azul","Marlim Azull","Marlim Azul",Replacer.ReplaceText,{"callerBranch.name"}),
#"Corrigindo OP Group" = Table.ReplaceValue(#"Corrigindo Marlim Azul1","TI - Infraestrutura","Infraestrutura",Replacer.ReplaceText,{"operatorGroup.name"}),
#"Corrigindo OP Group1" = Table.ReplaceValue(#"Corrigindo OP Group","TI - Sistemas","Sistemas",Replacer.ReplaceText,{"operatorGroup.name"}),
#"Corrigindo OP Group 2" = Table.ReplaceValue(#"Corrigindo OP Group1","N1 Sistemas","Sistemas",Replacer.ReplaceText,{"operatorGroup.name"}),
#"Corrigindo OP Group 3" = Table.ReplaceValue(#"Corrigindo OP Group 2","N1 Infraestrutura","Infraestrutura",Replacer.ReplaceText,{"operatorGroup.name"}),
#"Corrigindo OP Group 4" = Table.ReplaceValue(#"Corrigindo OP Group 3","N2 Sistemas","Sistemas",Replacer.ReplaceText,{"operatorGroup.name"}),
#"Corrigindo OP Group 5" = Table.ReplaceValue(#"Corrigindo OP Group 4","N2 Infraestrutura","Infraestrutura",Replacer.ReplaceText,{"operatorGroup.name"}),
#"Corrigindo P&D 2" = Table.ReplaceValue(#"Corrigindo OP Group 5","Pesquisa e DesenvolvimentoI","Pesquisa e Desenvolvimento",Replacer.ReplaceText,{"caller.department.name"}),
#"Corrigindo C&L2" = Table.ReplaceValue(#"Corrigindo P&D 2","Compras","Compras e Logística",Replacer.ReplaceValue,{"caller.department.name"}),
#"Corrigindo INSTCON" = Table.ReplaceValue(#"Corrigindo C&L2","Instrumentação","Instrumentação e Controle",Replacer.ReplaceValue,{"caller.department.name"}),
#"Changed Type" = Table.TransformColumnTypes(#"Corrigindo INSTCON",{{"targetDate", type datetime}, {"sla.targetDate", type datetime}, {"onHoldDate", type datetime}, {"callDate", type datetime}, {"creationDate", type datetime}, {"modificationDate", type datetime}, {"completedDate", type datetime}, {"closedDate", type datetime}}),
#"Expanded creator" = Table.ExpandRecordColumn(#"Changed Type", "creator", {"name"}, {"creator.name"}),
#"Duplicated Column1" = Table.AddColumn(#"Expanded creator", "requestT", each Text.Middle([request], Text.PositionOf([request], ":") + 1, Text.Length([request]))),
Custom1 = Table.AddColumn(#"Duplicated Column1", "requestTT", each Text.Middle([requestT], Text.PositionOf([requestT], ":") + 1, Text.Length([requestT]))),
#"Replaced Value" = Table.ReplaceValue(Custom1,"Descrição detalhada do problema:","",Replacer.ReplaceText,{"requestTT"}),
#"Removed Columns" = Table.RemoveColumns(#"Replaced Value",{"requestT"})
\end{lstlisting}
\begin{center}
	\text{Código da Transformação de Dados}
\end{center}
%\include{Capitulos/Cap4-Resultados}
%\include{Capitulos/Cap5-Conclusao}
% ----------------------------------------------------------
% -- Elementos Pós-Textuais: -------------------------------
\postextual  
\bibliography{bibliografia} % Referências bibliográficas
%\include{PosTextuais/Apendice}
%\include{PosTextuais/Anexo}
%\includepdf[pages={1,3-5}]{PosTextuais/includepdfpages.pdf}
% ----------------------------------------------------------
\end{document}  % fim do documento
